\appendix{Представление графического материала}

Графический материал, выполненный на отдельных листах,
изображен на рисунках А.1--А.\arabic{числоПлакатов}.
\setcounter{числоПлакатов}{0}

\renewcommand{\thefigure}{А.\arabic{figure}} % шаблон номера для плакатов

\begin{landscape}
	
	\begin{плакат}
		\includegraphics[width=0.82\linewidth]{VKR1.eps}
		\заголовок{Сведения о ВКРБ}
		\label{pl1:image}      
	\end{плакат}
	
	\begin{плакат}
		\includegraphics[width=0.82\linewidth]{VKR2.eps}
		\заголовок{Цель и задачи работы}
		\label{pl2:image}      
	\end{плакат}
	
	\begin{плакат}
		\includegraphics[width=0.82\linewidth]{VKR3.eps}
		\заголовок{Диаграмма прецедентов}
		\label{pl3:image}      
	\end{плакат}
	
	\begin{плакат}
		\includegraphics[width=0.82\linewidth]{VKR4.eps}
		\заголовок{Архитектура системы}
		\label{pl4:image}      
	\end{плакат}
	
	\begin{плакат}
		
\includegraphics[width=0.82\linewidth]{VKR5.eps}
		\заголовок{Архитектура MTCNN}
		\label{pl5:image}      
	\end{плакат}
	
	\begin{плакат}
		
\includegraphics[width=0.82\linewidth]{VKR6.eps}
		\заголовок{Архитектура InceptionResnetV1}
		\label{pl6:image}      
	\end{плакат}
	
	\begin{плакат}
		
\includegraphics[width=0.82\linewidth]{VKR7.eps}
		\заголовок{Архитектура ResNet18}
		\label{pl7:image}      
	\end{плакат}
	
	\begin{плакат}
		
\includegraphics[width=0.82\linewidth]{VKR8.eps}
		\заголовок{Архитектура сиамской нейронной сети}
		\label{pl8:image}      
	\end{плакат}
	
	\begin{плакат}
		
\includegraphics[width=0.82\linewidth]{VKR9.eps}
		\заголовок{Метод получения изображения отпечатков пальцев}
		\label{pl9:image}      
	\end{плакат}
	
	\begin{плакат}
		\includegraphics[width=0.82\linewidth]{VKR10.eps}
		\заголовок{Главные окна программы}
		\label{pl10:image}      
	\end{плакат}
	
	\begin{плакат}
		\includegraphics[width=0.82\linewidth]{VKR11.eps}
		\заголовок{Заключение}
		\label{pl11:image}      
	\end{плакат}
	
\end{landscape}
